% Ad Familiares 13.68 (ad-familiares-13-68) --- English translation
% Translated by Claude (Anthropic) from the Latin (Perseus).
% Style guide: STYLE.md.

\ciceroLetterOpener{M. TVLLIVS CICERO P. SERVILIO ISAVRICO PROCOS. CONLEGAE S. P.}

\ciceroSection{1} Your letter was very welcome to me, and from it I learned the course of your voyages: for in it you signalled that you remembered our bond, than which nothing could have given me more pleasure. For the rest, it will be far more welcome still if you will write to me, in a familiar vein, on the state of the commonwealth --- that is, on the condition of your province and on your arrangements there. Such things, though I shall hear of them from many on account of your renown, I will most gladly learn from your own letters.

\ciceroSection{2} For my part, I shall not often write to you what I think on the largest questions of the commonwealth, on account of the danger of letters of that sort; but what is being done, I shall write of more often. Still, I do seem to hope that Caesar, our colleague, will see to it --- and indeed does see to it --- that we have some sort of commonwealth at all; and to his deliberations it was of great moment that you should belong. But if it is more useful to you, that is to say more glorious, to be set over Asia and to safeguard that ill-affected portion of the commonwealth, then for me too the same course which serves your interest and your fame must be the more desirable.

\ciceroSection{3} As for me, whatever I judge to bear on your standing I shall attend to with the highest zeal and care; and above all I shall, with every mark of respect, look after that most illustrious man, your father --- which is what I am bound to do, both for the long-standing tie of our friendship, and for the kindnesses your family has shown me, and for his own dignity.
